% ============================================================================
% Дискретная квантовая механика (ДКМ): Формальные основания
% Академическая монография (book/sources/volumes/vol-formal-foundations-monograph.tex)
% ============================================================================

\chapter*{Предисловие. Предмет монографии и дедуктивное строение теории снизу вверх}
\markboth{Предисловие}{Предисловие}
\addcontentsline{toc}{chapter}{Предисловие. Предмет монографии и дедуктивное строение теории снизу вверх}

Настоящая монография не создавалась с предварительным намерением построить «единую теорию всего». Любая попытка умозрительно сочинить всеобъемлющую физическую схему поверх существующих теорий неизбежно страдает произволом в выборе аксиом и скрытой подгонкой параметров. Исторический и методологический генезис представленных здесь результатов был прямо противоположным и следовал классической традиции теоретической физики школы Л.\,Д.~Ландау: \emph{задать минимальные, физически неизбежные условия симметрии и сохранения и строго вывести тот единственный вид закона движения, который эти условия допускают}.

Тот факт, что при первой же записи дискретной циркуляции вокруг граней куба квантовая механика (деление разности фаз на непрерывный волновой остаток и дискретный корпускулярный перенос) и электродинамика Максвелла сомкнулись в одно и то же алгебраическое тождество, а затем из того же закона без введения каких-либо подгоночных коэффициентов последовательно вывелись константы электромагнитного, слабого, сильного и гравитационного взаимодействий, спектр атомных и ядерных оболочек, законы планетогенеза, молекулярной биофизики и макроскопическая механика, стал для автора полной неожиданностью и потребовал многократной независимой перепроверки. В настоящем томе, после того как вся эта цепочка прошла тотальную машинную верификацию в ядре доказательства теорем Lean~4 (\texttt{0 sorryAx}, \texttt{0 Classical.choice}), теория излагается в своём итоговом виде --- строго дедуктивно снизу вверх по всей лестнице масштабов.

\section*{§\,0.1. Метод построения: восхождение снизу вверх без постулатов}
\addcontentsline{toc}{section}{§\,0.1. Метод построения: восхождение снизу вверх без постулатов}

В классической традиции теоретической физики теория считается завершённой тогда, когда из минимальных физических условий сохранения и симметрии однозначно выводится закон движения и все его наблюдаемые следствия. В настоящем томе этот метод доведён до своего логического предела: прежде чем записывать уравнения движения, мы строго доказываем, почему у физической природы не может существовать никакой внешней памяти за её пределами, и почему неделимый элементарный объём пространства обязан описываться конечным замкнутым регистром состояний.

Из этого доказанного свойства развёртывается непрерывная восьмиступенчатая лестница следствий, составляющая восемь глав монографии:
\begin{itemize}[leftmargin=*,noitemsep]
\item \textbf{Глава~1 (Доказательство отсутствия внешней памяти у природы, вывод трёхмерности пространства, кубооктаэдра, динамики Верле и первого рождения вихря):} из замкнутости физической природы по взаимодействиям и когомологической аннигиляции внешнего ядра доказывается отсутствие внешней памяти у природы. Физическая неделимость элементарного объёма (дающая одно неделимое состояние собственной занятости центра) и плотное заполнение пространства без пустот вдоль координатных осей порождают состояния ветвления ячейки, взаимно-однозначная замкнутость которых с упорядоченными парами осевых переходов однозначно выделяет трёхмерное пространство, девятибитный однооборотный фазовый регистр, двухшаговый симплектический закон эволюции Верле и двенадцативершинную координационную оболочку кубооктаэдра. Далее прослеживается детерминированная траектория решётки от единственного девятибитного затравочного состояния на нулевом такте через киральный храповик доменной стенки к первому рождению квантового вихря и вихревых колец на втором такте, бариогенезису и первичному нуклеосинтезу, а также строится целочисленная теория механизмов и машин для степенных и винтовых каскадов.
\item \textbf{Глава~2 (Дискретная экспонента Эйлера, однородная проективная арифметика и выход в непрерывный анализ):} на циклическом фазовом кольце строится дискретная экспонента Эйлера с тождественно постоянным разностным вронскианом, вычисляются в замкнутом виде дискретные интегралы Френеля, Якоби, Дирихле, Парсеваля и Бесселя, вводится однородная проективная арифметика без деления в квадратичных и круговых числовых полях, строится замкнутая башня разностных производных частного произвольного порядка, выводятся проективные операторы логарифма, рациональной степени и размерности Хаусдорфа--Морана, и доказывается сходимость сумм по слоям кубооктаэдра к классическому интегралу Римана и кубатуре Симпсона.
\item \textbf{Глава~3 (Булева полнота Поста, редукция 3-SAT, конструктивные когомологии де~Рама и барьеры внутренних вычислителей):} доказывается функциональная полнота двух ортогональных каналов девятибитного сумматора по теореме Поста, точная полиномиальная редукция задачи выполнимости, конструктивные теоремы Стокса и Гаусса--Остроградского, конечно-аддитивная мера объёма без парадокса Банаха--Тарского, четыре физических барьера для любого конечного вычислителя внутри решётки и встроенный информационно-теоретический режим одноразового блокнота Вернама--Шеннона.
\item \textbf{Глава~4 (Дискретная и непрерывная квантовая механика, электродинамика Максвелла и граница применимости континуума):} деление разности фаз с остатком расщепляет динамику на непрерывный волновой остаток и дискретный корпускулярный перенос, детерминированно выводя вероятностное правило Борна, фундаментальный порог квантования (мёртвую зону затвора), подпороговый коридор обратимых вычислений без выделения тепла Ландауэра, скорость связи вдоль ребра кубооктаэдра, квантовую границу Цирельсона и полную систему законов теоретической электротехники.
\item \textbf{Глава~5 (Калибровочная группа Стандартной модели, спектр частиц, гравитационное преломление и космология):} симплициальный кликовый комплекс тринадцатиузлового ядра определяет явные матричные коммутаторы алгебры Ли калибровочной группы Стандартной модели, четырёхпетлевое разложение по степеням простого гауссова кулоновского модуля даёт замкнутую алгебраическую форму постоянной тонкой структуры без участия трансцендентного числа пи, геометрия тринадцатиузлового одетого ротора электрона и барионных треугольных граней выводит спектр масс лептонов и кварков, углы смешивания нейтрино и кварков, магнитные моменты нуклонов и энергию связи атома водорода, а преломление на фазовом импедансе ворот устраняет гравитационную сингулярность, описывает фазы коры нейтронных звёзд и объясняет плоские кривые вращения галактик за счёт вакуумного гало Вороного без постулирования частиц тёмной материи.
\item \textbf{Глава~6 (Атомные и ядерные оболочки кубооктаэдра, планетогенез по шести зонам диска, сверхпроводимость и когерентный коллапс барьера):} радиальные координационные сферы кубооктаэдра порождают периодическую систему элементов до восьмого периода, четыре физических предела устойчивости ядер, двухсферный перенос урана на остров стабильности сверхтяжёлых элементов, полный закон планетогенеза по всем шести зонам протопланетного диска, четыре мультипликативных тождества Безу, правила орбитальной симметрии химической связи, сверхпроводящее спаривание Купера и когерентное снижение кулоновского барьера при синфазной фокусировке двенадцати лучей.
\item \textbf{Глава~7 (Молекулярная биофизика, планетарные фотонные окна, биометаллы, генетический код ДНК, фолдинг белков и морфогенез):} из делителя водородной связи выводятся три планетарных фотонных окна биосферы, автономно открываются шесть кофакторных биометаллов металлоферментов в таблице Менделеева, доказывается защита всех четырёхсот парных стыков канонических кодонов ДНК от внефазного старт-кодона, ко-трансляционная укладка белка за линейное время, сквозная компиляция топологии органа в последовательность ДНК, три режима метаболической аллометрии, инварианты шестислойной кортикальной колонки неокортекса и межэтажный изоморфизм замкнутых оболочек от квантового до биологического этажа.
\item \textbf{Глава~8 (Бесфононный кремниевый фотокаскад, общая конструкция предела по счётным моноидам стадий и гёделевская неизбежность лабораторного арбитража):} выводится резонансное условие бесфононного синглетного деления синего фотона в четырёхслойном кремниевом фотокаскаде и архитектура безветвевого кремниевого вычислительного стека, строится общий оператор предела по счётным моноидам стадий и проективным башням гомоморфизмов алгебраических систем А.\,И.~Мальцева, доказывается двусторонняя эквивалентность дискретного и непрерывного носителей, и по диагональной теореме Ловера--Гёделя--Тарского обосновывается необходимость прямого лабораторного эксперимента как единственного внешнего семантического арбитра теории.
\end{itemize}

\section*{§\,0.2. Четыре методологических опоры монографии}
\addcontentsline{toc}{section}{§\,0.2. Четыре методологических опоры монографии}

Чтобы результаты такого масштаба могли быть подвергнуты строгому, непредвзятому научному аудиту, изложение в настоящем томе опирается на четыре принципа:

\begin{enumerate}[leftmargin=*]
\item \textbf{Единство физических механизмов природы.}
  Природа не выбирает, по какой ведомственной кафедре --- квантовой электродинамики, ядерной физики, физики твёрдого тела, химии, планетологии или молекулярной биологии --- ей работать в данный момент. Она либо реализует единый закон обновления на дискретном носителе, либо нет. Если закон верен на уровне элементарной ячейки, то одни и те же алгебраические инварианты координационной оболочки кубооктаэдра и девятибитного фазового кольца обязаны проявляться и в постоянной тонкой структуры, и в спектрах ядерных и электронных оболочек, и в металлоферментах живой клетки, и в макроскопической электродинамике.

\item \textbf{Тотальная машинная верификация в ядре Lean~4 (\texttt{0 sorryAx}, \texttt{0 Classical.choice}).}
  В теоретической физике главным источником ошибок всегда был человеческий фактор: неявные допущения в длинных аналитических выкладках, округления и некорректные предельные переходы. В настоящей работе этот риск устранён радикально: вся цепочка определений и теорем формализована на языке зависимых типов \textbf{Lean~4} (\texttt{lean/MtCarrier/} и \texttt{papers/}). Во всех операторах размерность пространства передаётся как свободный натуральный параметр, а каждое утверждение доказано конструктивно без единого пропуска и без использования неконструктивной аксиомы выбора. Любой исследователь может независимо проверить все доказательства командой \texttt{lake build}.

\item \textbf{Строгая преемственность и принцип соответствия Бора.}
  Предлагаемая дискретная теория \textbf{не опровергает и не отменяет} ни классическую механику, ни уравнения Максвелла, ни общую теорию относительности, ни стандартную непрерывную квантовую механику Шрёдингера--Дирака, ни классический математический анализ Римана--Лиувилля. Напротив, в Главах~2--4 строго доказывается, что при большом числе узлов и тактов суммы по слоям кубооктаэдра и кусочно-линейные траектории на фазовом кольце сходятся со вторым порядком точности к классическим интегралам Римана и дифференциальным уравнениям непрерывной квантовой механики на континуальном носителе, а квадратичное детектирование на фазовом кольце тождественно воспроизводит правило Борна.
  Однако у континуальной аппроксимации появляется \textbf{точная физическая граница применимости}: там, где непрерывный формализм при стремлении шага к нулю порождает ультрафиолетовые расходимости, гравитационные сингулярности деления на ноль, парадокс Банаха--Тарского и тепловой предел стирания информации Ландауэра, дискретный носитель остаётся конечным, побитово обратимым и свободным от расходимостей.

\item \textbf{Инженерная реализуемость всего, что не запрещено законами природы.}
  Технологические следствия теории --- от обратимых безветвевых вычислений без выделения тепла Ландауэра и когерентного снижения кулоновских барьеров катализа до бесфононных каскадов кремниевой фотовольтаики --- представляют собой задачу \textbf{точного согласования фазовых и оболочечных импедансов} вместо традиционного несогласованного порогового преодоления энергетических барьеров.
\end{enumerate}

\section*{§\,0.3. Двусортное проективно-ретракционное соответствие оснований естественных наук: Трёхзаконная целочисленная механика Ньютона}
\addcontentsline{toc}{section}{§\,0.3. Двусортное проективно-ретракционное соответствие оснований естественных наук}

Центральный онтологический вывод настоящей монографии состоит в том, что \textbf{все разделы фундаментальной физики, химии, молекулярной биологии и инженерии сводятся через двусортное проективно-ретракционное соответствие дискретного и непрерывного носителей к классической механике трёх законов Ньютона в точных целых числах на решётке кубооктаэдра}:
\begin{enumerate}[label=\textbf{Закон Ньютона \Roman*.},leftmargin=*,noitemsep]
\item \textbf{Закон инерции в подпороговом беспотерьном коридоре:}
  при отсутствии надпорогового вращающего момента вторая разность фазы равна нулю, а первая разность сохраняется постоянной. Стационарные атомные орбиты, сверхпроводящие токи, баллистические экситоны и обратимые логические вентили движутся по инерции без излучения и без выделения тепла Ландауэра.
\item \textbf{Симплектический закон динамики второго порядка с переменной инерцией узла:}
  единственная обратимая во времени двухшаговая рекурсия Верле представляет собой второй закон Ньютона на фазовом кольце, где локальный фазовый импеданс узла играет роль переменного момента инерции, ограниченного снизу вакуумным полом и сверху потолком противофазного стазиса.
\item \textbf{Равенство действия и противодействия по всем двенадцати лучам кубооктаэдра:}
  в силу антисимметрии нечётной проекции парные потоки между любыми двумя соседними узлами в сумме тождественно равны нулю, напрямую обеспечивая точное сохранение импульса, момента импульса, электрического заряда и вихревого потока.
\end{enumerate}

Все эффекты, для которых в континуальных теориях XX~века вводились отдельные постулаты, возникают не из отказа от трёх законов Ньютона, а из пяти механических свойств самой дискретной решётки кубооктаэдра, сопоставленных в Таблице~\ref{tab:mono-newtonian-homeomorphism}.

\begin{table}[htbp]
\centering
\caption{Двусортное проективно-ретракционное соответствие носителей: сведение континуальных физических и инженерных постулатов к трёхзаконной целочисленной механике Ньютона.}
\label{tab:mono-newtonian-homeomorphism}
\footnotesize
\begin{tabular}{@{}p{0.21\textwidth}p{0.35\textwidth}p{0.37\textwidth}@{}}
\toprule
\textbf{Раздел и свойство решётки} & \textbf{Континуальная идеализация на непрерывном носителе} & \textbf{Дискретная ньютоновская механика на решётке кубооктаэдра} \\
\midrule
\textbf{1. Кинематика и релятивизм}\newline (Предел одного ребра за такт) &
Постулирует сокращение непрерывного пространства-времени и трактует квантовую запутанность как нелокальное дальнодействие. &
Жёсткая решётка и дискретные такты; предельная скорость связи вдоль диагонального ребра кубооктаэдра замедляет внутренние роторы движущихся частиц и даёт квантовую границу корреляций Цирельсона. \\
\addlinespace[3pt]
\textbf{2. Гравитация и космология}\newline (Импеданс ворот и ячейка Вороного) &
Искривляет пустой континуум, порождая сингулярности деления на ноль в чёрных дырах и требуя постулирования частиц тёмной материи. &
Второй закон Ньютона в среде с переменным фазовым импедансом; потолок противофазного стазиса исключает сингулярности, а геометрия ячейки Вороного даёт плоские кривые вращения галактик без тёмной материи. \\
\addlinespace[3pt]
\textbf{3. Квантовая механика}\newline (Порог зацепления Гейзенберга) &
Заменяет траектории волнами вероятности в гильбертовом пространстве, постулирует коллапс при измерении и тепловой предел стирания бита. &
Детерминированная траектория Верле в каждой ячейке; целочисленный порог квантования даёт инерционное движение без излучения в подпороговом коридоре и выводит правило Борна из статистики переполнения кольца. \\
\addlinespace[3pt]
\textbf{4. Электродинамика и теория поля}\newline (Фазовое кольцо и противофазный потолок) &
Постулирует калибровочные расслоения и требует перенормировки расходящихся ультрафиолетовых петлевых интегралов. &
Деление разности фаз с остатком на период кольца даёт квантованный вихревой поток и уравнения Максвелла, а полупериодный противофазный потолок ограничивает число излучаемых квантов за такт. \\
\addlinespace[3pt]
\textbf{5. Химия, биология и инженерия}\newline (Оболочки и седло четверти оборота) &
Опирается на эмпирические правила заполнения орбиталей, экспоненциальный перебор конформаций белка и однопереходный предел фотоэлементов. &
Слои и клики кубооктаэдра порождают таблицу Менделеева, зоны планетогенеза и биометаллы; седло четверти оборота объясняет когерентный катализ, линейный по времени фолдинг белка и бесфононный кремниевый фотокаскад. \\
\bottomrule
\end{tabular}
\end{table}

